\cotefile{93}
\scanpage{17}\mspage{- 14 -}
\makeatletter\tagsleft@true\makeatother

$2$ nilpotent, ce qui permet de mettre la courbe sous la forme (4.1)
\[
y^2+xy=x^3+a_2x+a_6
\]
et un isomorphisme d'une telle courbe sur une autre est décrit par $s$ vérifiant (4.3) pour $u=1+2s$. Lorsque $j=j'$, la première de ces équations
\[
(1+2s)^2\,a'_2=a_2-s-s^2
\]
définit un revêtement étale de degré deux, et elle implique la seconde ; pour le voir, on se ramène à supposer $a_2=a'_2$ et on remarque que d'après (4.4)
\[
1/j=a_6+2\,R(a_2,a_6)
\]
pour une fraction rationnelle $R$.

Pour vérifier (II) dans les cas non couverts par (III), on peut supposer que $S$ est le spectre d'un corps algébriquement clos et que $c_4=0$ ou $c_6=0$ ; i.e. que $j=0$ ou $j=1728$. Soit $p$ la caractéristique du corps,

a) $p\ne 2,\,3$ : on met la courbe sous la forme (2.1) ; si $c_4=0$ (resp. $c_6=0$) on aura $a_4=0$ (resp. $a_6=0$) et on peut par homothétie transformer la courbe en la courbe type
\[
y^2=x^3+1\qquad(\text{si }c_4=0)\tag*{(6.5)}
\]
\[
y^2=x^3-x\qquad(\text{si }c_6=0)\tag*{(6.6)}
\]

b) $p=3$ : on met la courbe sour la forme (6.1) ; résolvant (6.2) pour $a'_4=1$ et $a'_6=0$, on met la courbe sous la forme type

% archive sheet number: 16
\scanpage{18}

\uline{\unsure{question|Question} connexe} : soit $E$ ell sur $k$ parfait.
\[
M^{\wedge}=\text{sch formel \unsure{mod} de }E\;(\text{sur }W(k))
\]
\[
\boxed{\operatorname{Spec} W(k)[j]^{\wedge}\sim M^{\wedge}/\operatorname{Aut}(E)}
\quad\longrightarrow\quad\text{étude }\infty^{\mathrm{ale}}\text{ à faire de }j.
\]

\diagram{\(\begin{array}{llll}
\text{a)}&j=0,\ p\ne2,3\ :&j\sim t^3&t\ \text{paramètre régulier relatif}\\
\text{b)}&j=12^3,\ p\ne2,3&j\sim t^2&\text{''}\\
\text{c)}&\text{\uline{$j=0,\ p=3$}}\ :&\multicolumn{2}{l}{\text{groupe }6\text{ éléments }-}\\
\text{d)}&\text{\uline{$j=0,\ p=2$}}\ :&\multicolumn{2}{l}{\text{''}\quad12\quad\text{''}\qquad-}
\end{array}\)}
\margin{$\begin{aligned}
\unsure{\parallel|11} &=\mathrm{car}\ 3:\quad 6\mid v(j)\\
&=\mathrm{car}\ 2:\quad 12\mid v(j)
\end{aligned}$}

\unsure{$9$|$\rho$} ? semi direct de $\unsure{\mathbb{Z}|Z}/3$ par $\{\pm1\ \pm i\ \pm j\ \pm k\}$ \unsure{ou|où} $\mathbb{Z}/3$ agit par $(i,j,k)$ \quad $-$

% archive sheet number: 17
\scanpage{19}\mspage{- 15 -}

\[
y^2=x^3-x\qquad(\text{si }c_4=c_6=0)\tag*{(6.7)}
\]

c) $p=2$ : ici, on met la courbe sour la forme (6.3) et on résoud (6.4) avec $a_3=1$, $a_4=a_6=0$ pour obtenir la forme type
\[
y^2+y=x^3\qquad(\text{si }c_4=c_6=0)\ .\tag*{(6.8)}
\]

Pour compléter le tableau, reste à décrire les automorphismes des courbes elliptiques.

\uline{Proposition} 6.9. \uline{Soit} $E$ \uline{une courbe de genre un sur une base} $S$, \uline{sans fibre additive}

(I) \uline{Si} $c_4$ et $c_6$ \uline{sont inversibles,} \uline{les seuls automorphismes de} $E$ \uline{sont l'identité et la symétrie}.
\margin{\handwritten{[vrai seulement en \unsure{faisceautisant|fais-ceautisant} !]}}

(II) \uline{Si} $2$ et $3$ \uline{sont inversibles}, \uline{l'application canonique}
\[
\underline{\operatorname{Aut}}(E)\handwritten{\longrightarrow}G_m\tag*{(6.6.)}
\]
\uline{donnée par l'action sur les différentielles invariantes}, \uline{est une immersion fermée}.

(III) \uline{Si} $2$ \uline{et} $3$ \uline{sont inversibles}, \uline{et} $c_4=0$ (resp. $c_6=0$) \uline{l'image de} (6.6) \uline{est} $\mu_6$ (resp. $\mu_4$).

\uline{Supposons maintenant que} $S$ \uline{soit le spectre d'\unsure{ur|un} corps}

(IV) \uline{Si} $3=c_4=c_6=0$, $\operatorname{Aut}(E)$ \uline{est géométriquement \unsure{isomorphe} au produit semi-direct} (\uline{non commutatif}) \uline{de} $\mathbb{Z}/(4)$ \uline{par} $\mathbb{Z}/(3)$, \uline{et l'image de} (6.6) \uline{est} $\mu_4$.

(V) \uline{Si} $2=c_4=c_6=0$, $\operatorname{Aut}(E)$ \uline{est géométriquement d'ordre} $24$, \uline{et l'image de} (6.6) \uline{est} $\mu_3$.
\handwritten{\unsure{$9$|$\rho$|$?$}}

% archive sheet number: 18
\scanpage{20}\mspage{- 16 -}

\uline{Preuve} : (I) D'après (6.3), il n'y a en effet pas plus de deux automorphismes.

(II) et (III) Lorsque les courbes sont mises sous la forme (2.1), il résulte de (2.5) que les automorphismes sont décrits parleur action sur $\pi$ et que $\lambda\in G_m$ défini un automorphisme si et seulement si
\[
c_4=\lambda^4c_4\qquad\text{et}\qquad c_6=\lambda^6c_6\quad,
\]
d'où les assertions.

(IV) Si la courbe est mise sous forme (6.1), les automorphismes sont donnés par (1.2) : $x'\handwritten{\rightsquigarrow}x$, pour $s=t=0$ et
\[
\handwritten{\Bigg\{}\begin{aligned}
&u^4=1\\
&r^3+ra_4+a_6(1-u^2)=0
\end{aligned}
\]
et la loi de composition par
\[
(u',r')\quad(u,r)=(u\,u',\,u^{\prime 2}r+r')
\]

Lorsque la courbe est mise sous forme (6.7), ces équations se réduisent à $u\in\mu_4$ et $r\in F_3$.

(V) Si la courbe est mise sous forme (6.3), les automorphismes sont donnés par (1.2) : $x'\handwritten{\rightsquigarrow}x$, pour $r=s^2$ et
\[
\handwritten{\Bigg\{}\begin{aligned}
&u^3=1\\
&s^4+sa_3+a_4(1-u)=0\\
&t^2+ta_3+s^6+s^2a_4=0
\end{aligned}
\]

% archive sheet number: 19
\scanpage{21}

\uline{Cor} : pour $\unsure{S|B}$ plat sur $\mathbb{Z}$, $\Gamma(S,\unsure{\mathcal{O}|O}_S)[\,\ldots\,]$

[ou seul\textsuperscript{\uline{\unsure{mt|nt}}} $2$ et $3$ non div. de $0$]

% archive sheet number: 20
\scanpage{22}\mspage{- 17 -}

et la loi de composition par
\[
(u',s',t')\mathbin{\unsure{o|\circ}}(u,s,t)=(u\,u',\,u's+t',\,u^{\prime 3}t+t'+u^{\prime 2}rs')\ .
\]

Lorsque la courbe est mise sous forme (6.8), ces équations se réduisent à
\[
\begin{aligned}
u\in\mu_3\qquad s\in\mathbb{F}_4\qquad t^2+t=s^3\quad &(=0\text{ ou }1\text{ selon que}\\
&s=0\text{ ou }s\ne0)\ .
\end{aligned}
\]

7. \uline{Anneau des formes modulaires entières}.

On appelle forme modulaire entière de poids $n$ une loi qui à chaque courbe de genre un sur une base $S$ associe une section de $\omega^{\otimes n}$, et ce de façon compatible au changement de base. Il reviendrait au même de se limiter aux courbes sans fibres additives.

\uline{Proposition} 7.1. \uline{L'anneau des formes modulaires entières est engendré sur} $\mathbb{Z}$ \uline{par} $c_4$, $c_6$ \uline{et} $\Delta$ \uline{soumis à la seule relation}
\[
c_4^3-c_6^2=1728\,\Delta\quad.\tag*{(7.2.)}
\]

Appliquant la définition à la courbe (1.3), on voit qu'une forme modulaire \unsure{\.{e}ntière|éntière|entière} de poids $n$ s'exprime comme polynôme de\margin{\handwritten{$\unsure{\sigma|6|o}$ ?}} poids $n$ en les $a_i$. Pour $2$ et $3$ inversibles, il résulte de 2.5 que ces formes sont les polynômes en $c_4$ et $c_6$, et en particulier que $c_4$, $c_6$ et $\Delta$ ne satisfont pas d'autres relations que (7.2). Reste à prouver que si un polynôme en $c_4$, $c_6$ et $\Delta$ est identiquement nul en caractéristique $2$ (resp. $3$), alors il est divisible par $2$ (resp.$3$) dans $\mathbb{Z}[c_4,c_6,\Delta]/c_4^3-c_6^2-1728\,\Delta$, ce qui signifie

% archive sheet number: 21
\scanpage{23}

$a_1$ et $\pm b_2$ sont les invariants de \unsure{Hasse|Hase} Witt.

% archive sheet number: 22
\scanpage{24}\mspage{- 18 -}

qu'il est divisible par $c_4^3-c_6^2$ dans $\mathbb{Z}/(2)[c_4,c_6,\Delta]$ (resp. $\mathbb{Z}/(3)[c_4,c_6,\Delta]$). Cela résulte de la structure de l'anneau des formes modulaires de caractéristique $2$ (resp. $3$).

\uline{Proposition} 7.2. (I) \uline{L'anneau des formes modulaires de caractéristique} $2$ \uline{est} $\mathbb{Z}/(2)[a_1,\Delta]$, \uline{et} $c_4=a_1^4$, $c_6=a_1^6$

(II) \uline{L'anneau des formes modulaires de caractéristique} $3$ \uline{est} $\mathbb{Z}/(3)[b_2,\Delta]$, et $c_4=b_2^2$, $c_6=-b_2^3$.

Tout d'abord, il résulte de $(1.6)_1$ et $(1.7)_2$ que $a_1$ (resp. $b_2$) est une forme modulaire de caractéristique $2$ (resp. $3$).

En caractéristique $2$, cherchons tout d'abord les formes modulaires relatives aux seules courbes telles que $a_1$ soit inversible. Ces courbes se mettent sous forme (cf. (4.1))
\[
y^2+a_1xy=x^3+a_2x^2+a_6\qquad\quad\pi=\frac{dx}{a_1x}
\]
et, compte tenu de (4.3), on voit que ces formes s'écrivent
\[
\begin{aligned}
&a_1^{-n}\,P(a_1,a_6)\qquad\text{pour }P\text{ isobare, ou encore}\\
&a_1^{-m}\,Q(a_1,\Delta)
\end{aligned}
\]
et cela garde un sens pour $a_1=0$, $\Delta\ne0$ si et seulement si on avait en fait affaire à un polynome en $a_1$ et $\Delta$.

En caractéristique $3$, raisonnant de même d'abord sur les courbes telles que $c_4$ soit inversible, et remarquant que $d_6=0$ équivaut alors à $a_4=0$ dans (3.1), on voit que toute forme modu%
% archive sheet number: 23
\scanpage{25}\mspage{- 19 -}%
laire s'écrit
\[
b_2^{-n}\,P(b_2,\Delta)\qquad\text{pour }P\text{ isobare}\quad,
\]
et que cette forme modulaire est définie aussi pour $j=0$ si et seulement si elle s'écrit comme polynôme en $b_2$ et $\Delta$.

% archive sheet number: 24
\scanpage{27}

Orientation

% archive sheet number: 26
\scanpage{28}
\textbf{C. R. Acad. Sc. Paris, t. 292 (19 janvier 1981)\hfill Série I --- }\mspage{213}

TOPOLOGIE. --- \textit{Théorie axiomatique de l'orientation des variétés topologiques à bord.} Note (*) de \textbf{Marcus Wanderley}, présentée par Henri Cartan,

On définit ce qu'on entend par une théorie de l'orientation des variétés topologiques (à bord) de dimension $\leqq\mathrm{N}$, et on prouve l'unicité (à isomorphisme près) d'une telle théorie lorsque $\mathrm{N}\ne4$, en s'appuyant sur un résultat de Kirby (pour $n\ne4$, tout homéomorphisme de $\mathbf{R}^n$ de degré $+1$ est « stable »).

\textit{One defines what is understood by an orientation theory for the topological manifolds (with boundary) of dimension $\leqq\mathrm{N}$. One proves that such a theory is unique (up to isomorphism) if $\mathrm{N}\ne4$, using a result of Kirby (for $n\ne4$ every homeomorphism of $\mathbf{R}^n$ of degree $+1$ is ``stable'').}

\textsc{Notations}. --- Dans cette Note, $\mathrm{N}$ désigne un entier positif ou le symbole $\infty$. On désigne par $\mathscr{V}^{\bullet}_{\leqq\mathrm{N}}$ la catégorie dont les objets sont les variétés topologiques à bord de dimension $\leqq\mathrm{N}$, les morphismes étant les plongements ouverts. La sous-catégorie pleine de $\mathscr{V}^{\bullet}_{\leqq\mathrm{N}}$ formée des variétés topologiques (i. e. à bord vide) de dimension $\leqq\mathrm{N}$ est notée par $\mathscr{V}_{\leqq\mathrm{N}}$. Si $m$ est un entier positif $\leqq\mathrm{N}$, on désigne par $\mathscr{V}^{\bullet}_m$ (resp. $\mathscr{V}_m$) la sous-catégorie pleine de $\mathscr{V}^{\bullet}_{\leqq\mathrm{N}}$ (resp. $\mathscr{V}_{\leqq\mathrm{N}}$) des variétés pures de dimension $m$. Lorsque l'on a un foncteur contravariant :
\[
\mathrm{F}:\quad\mathscr{V}^{\unsure{0|\mathrm{o}}}_{\leqq\mathrm{N}}\to\mathrm{Ens},
\]
on le prolonge à un foncteur :
\[
\mathrm{F}^{\bullet}:\quad\mathscr{V}^{\bullet\;0}_{\leqq\mathrm{N}}\to\mathrm{Ens},
\]
par la formule :
\[
\mathrm{F}^{\bullet}(\mathrm{V})=\mathrm{F}(\operatorname{Int}(\mathrm{V}))\qquad [\mathrm{V}\in\operatorname{ob}(\mathscr{V}^{\bullet}_{\leqq\mathrm{N}})].
\]

\textsc{Définition}. --- Une théorie $\mathrm{T}=(\Omega,\delta,\nu)$ de l'orientation des variétés topologiques de dimension $\leqq\mathrm{N}$ consiste en les données (1), (2) et (3) ci-dessous, astreintes aux axiomes $(a)$, $(b)$, $(c)$ et $(d)$ ci-dessous.

\textsc{Données}. --- (1) (\textit{Foncteur}). On se donne un foncteur contravariant :
\[
\Omega:\quad\mathscr{V}^{0}_{\leqq\mathrm{N}}\to\mathrm{Ens},
\]
appelé \textit{foncteur orientation}. Le prolongement $\Omega^{\bullet}$ de $\Omega$ à $\mathscr{V}^{\bullet\;0}_{\leqq\mathrm{N}}$ est appelé également \textit{foncteur orientation}. Si $\mathrm{V}$ est un objet de $\mathscr{V}^{\bullet}_{\leqq\mathrm{N}}$ alors $\Omega^{\bullet}(\mathrm{V})$ est appelé \textit{l'ensemble des orientations de} $\mathrm{V}$ (au sens de $\mathrm{T}$).

(2) (\textit{Stokes}). On se donne pour tout objet $\mathrm{V}$ de $\mathscr{V}^{\bullet}_{\leqq\mathrm{N}}$ une application :
\[
\delta_{\mathrm{V}}:\quad\Omega^{\bullet}(\mathrm{V})\to\Omega\,(\partial\mathrm{V}),
\]
fonctorielle en $\mathrm{V}$ pour $\mathrm{V}$ variable, appelée \textit{application de Stokes}. Si $\omega$ est une orientation de $\mathrm{V}$ alors $\delta_{\mathrm{V}}(\omega)$ est appelée \textit{orientation unduite sur le bord} $\partial\mathrm{V}$ \textit{de} $\mathrm{V}$.

(3) (\textit{Normalisation}). On se donne pour la variété $\mathrm{P}$ réduite à un point une bijection :
\[
\nu:\quad\{\pm1\}\overset{\sim}{\to}\Omega(\mathrm{P}),
\]
appelée \textit{application de normalisation}.

\textsc{Axiomes}. --- $(a)$ (\textit{Exactitude}). $\Omega$ est un « faisceau », c'est-à-dire pour tout $\mathrm{V}$ objet de $\mathscr{V}_{\leqq\mathrm{N}}$ la \textit{restriction} $\Omega_{\mathrm{V}}$ de $\Omega$ à la catégorie des ouverts de $\mathrm{V}$ avec comme morphismes les inclusions, est un faisceau (\textit{cf.} [1]).

% archive sheet number: 27
\scanpage{29}\mspage{214}
\textbf{--- Série I\hfill C. R. Acad. Sc. Paris, t. 292 (19 janvier 1981)}

$(b)$ (\textit{Séparation}). $\Omega$ est un faisceau « (topologiquement) séparé », c'est-à-dire pour toute variété $\mathrm{V}$ $(\dim\mathrm{V}\leqq\mathrm{N})$ et tout couple $(\omega,\ \omega')$ d'orientations de $\mathrm{V}$ l'ensemble ouvert $\mathrm{C}_{\omega,\omega'}=\{x\in\mathrm{V}\mid\omega(x)=\omega'(x)\}$ est aussi fermé.

$(c)$ (\textit{Isomorphisme local}). Pour toute variété à bord $\mathrm{V}$ $(\dim\mathrm{V}\leqq\mathrm{N})$ l'application de Stokes $\delta_{\mathrm{V}}$ est un « isomorphisme d'espaces étalés sur $\mathrm{V}$ et $\partial\mathrm{V}$ », c'est-à-dire les conditions ci-dessous sont vérifiées :

(i) Si $\omega,\ \omega'\in\Omega^{\bullet}(\mathrm{V})$ « induisent » la même orientation sur le bord $\partial\mathrm{V}$ de $\mathrm{V}$ [i. e. $\delta_{\mathrm{V}}(\omega)=\delta_{\mathrm{V}}(\omega')$], alors $\omega$ et $\omega'$ coïncident sur un voisinage ouvert de $\partial\mathrm{V}$.

(ii) Si $\bar{\omega}\in\Omega(\partial\mathrm{V})$, alors $\bar{\omega}$ est « localement induite », c'est-à-dire pour tout $x\in\partial\mathrm{V}$ il existe un voisinage ouvert $\mathrm{U}$ de $x$ (en $\mathrm{V}$) et une orientation $\omega\in\Omega^{\bullet}(\mathrm{U})$ tel que $\delta_{\mathrm{V}}(\omega)=\bar{\omega}\mid\partial\mathrm{U}$ ($\bar{\omega}\mid\partial\mathrm{U}$ désigne la restriction de $\bar{\omega}$ à $\partial\mathrm{U}$).

On voit aussitôt que pour toute variété à bord $\mathrm{V}$ la conjonction des conditions (i) et (ii) ci-dessus équivaut à dire que l'application :
\[
\Omega^{\bullet}_{\mathrm{V}}\mid\partial\mathrm{V}\overset{\sim}{\to}\Omega_{\partial\mathrm{V}},
\]
déduite de l'application de Stokes est un \textit{isomorphisme d'espaces étalés de base} $\partial\mathrm{V}$.

Avec les données et axiomes précédents on démontre la proposition suivante :

\textsc{Proposition}. --- \textit{Pour tout objet} $\mathrm{V}$ \textit{de} $\mathscr{V}^{\bullet}_{\leqq\mathrm{N}}$ \textit{l'espace étalé} $\Omega_{\mathrm{V}}$ \textit{sur} $\mathrm{V}$ \textit{défini par} $\Omega$ \textit{est un revêtement à deux feuillets} (\textit{appelé} revêtement des orientations de $\mathrm{V}$).

Il résulte de cette proposition que le groupe $\mu_2=\{\pm1\}$ opère sur le faisceau $\Omega$. Si $\omega$ est une orientation d'une variété $\mathrm{V}$ on pose $-\omega=-1\,.\,\omega$; l'orientation $-\omega$ est appelée l'\textit{orientation opposée à} $\omega$.

$(d)$ (\textit{Anti-commutativité}). --- Pour tout $\mathrm{V}$ objet de $\mathscr{V}_{\leqq\mathrm{N}}$ et pour des sous-variétés $\mathrm{V}'$, $\mathrm{V}''$ de $\mathrm{V}$ tels que $\mathrm{V}'\cup\mathrm{V}''=\mathrm{V}$ et $\mathrm{V}'\cap\mathrm{V}''=\partial\mathrm{V}'=\partial\mathrm{V}''$, le diagramme :
\diagram{\begin{tikzcd}[ampersand replacement=\&,row sep=small,column sep=large]
\&\Omega^{\bullet}(\mathrm{V}')\arrow[dr,"{\delta_{\mathrm{V}'}}"]\&\\
\Omega(\mathrm{V})\arrow[ur,"{\unsure{r|r'}}"]\arrow[dr,"{r''}"']\&\&\Omega(\partial\mathrm{V}')=\Omega(\partial\mathrm{V}''),\\
\&\Omega^{\bullet}(\mathrm{V}'')\arrow[ur,"{\delta_{\mathrm{V}''}}"']\&
\end{tikzcd}}
\textit{où} $r'$ et $r''$ sont les restrictions, est « anti-commutatif », c'est-à-dire pour toute orientation $\omega$ de $\mathrm{V}$ on a la relation :
\[
\delta_{\mathrm{V}'}(\omega\mid\operatorname{Int}(\mathrm{V}'))=-\delta_{\mathrm{V}''}(\omega\mid\operatorname{Int}(\mathrm{V}'')).
\]

\textsc{Définition}. --- Soient $\mathrm{T}=(\Omega,\delta,\nu)$ et $\mathrm{T}'=(\Omega',\delta',\nu')$ deux théories de l'orientation des variétés topologiques en dimension $\leqq\mathrm{N}$. Un morphisme de $\mathrm{T}$ en $\mathrm{T}'$ est un morphisme fonctoriel de $\Omega$ en $\Omega'$ qui commute aux données de Stokes et de normalisation, en un sens évident.

\textsc{Théorème} 1 (\textit{cf.} [3]). --- $(a)$ \textit{Il existe une théorie de l'orientation des variétés topologiques en dimension} $\leqq N$.

$(b)$ \textit{Il existe au plus} un morphisme \textit{entre deux théories de l'orientation des variétés topologiques en dimension} $\leqq\mathrm{N}$ \textit{et un tel morphisme est un} isomorphisme.

$(c)$ \textit{Si} $\mathrm{N}\ne4$ \textit{il existe un} (unique) \textit{isomorphisme entre deux théories de l'orientation des variétés topologiques en dimension} $\leqq\mathrm{N}$.

On démontre ce théorème en utilisant les résultats ci-après.

% archive sheet number: 28
\scanpage{30}
\textbf{C. R. Acad. Sc. Paris, t. 292 (19 janvier 1981)\hfill Série I --- }\mspage{215}

\textsc{Théorème} 2 (cf. [3]). --- \textit{Soit} $m\in\mathbf{N}^{*}$ \textit{et soit} $\operatorname{Mon}(\mathbf{R}^m)$ \textit{le monoïde des applications continues et injectives de} $\mathbf{R}^m$ \textit{en} $\mathbf{R}^m$. \textit{Si} $\Omega$ \textit{est un faisceau sur} $\mathscr{V}_m$ [\textit{au sens de l'axiome} $(a)$] \textit{les conditions ci-dessous sont équivalentes} :

(i) $\Omega$ \textit{est localement constant} (cf. [4]);

(ii) $\Omega$ \textit{est séparé} [\textit{au sens de l'axiome} $(b)$];

(iii) \textit{Le monoïde} $\operatorname{Mon}(\mathbf{R}^m)$ \textit{opère sur} $\Omega(\mathbf{R}^m)$ \textit{par automorphismes}.

Soit $\mathrm{L.\ C.}\,(\mathscr{V}_m)$ la catégorie des faisceaux localement constants sur $\mathscr{V}_m$ et soit $\Pi(m)$ le groupe symétrisé du monoïde $\operatorname{Mon}(\mathbf{R}^m)$.

\textsc{Théorème} 3 (cf. [3]). --- \textit{Le foncteur} [cf. \textit{théorème} 2, (iii)] :
\[
\Omega\to\Omega(\mathbf{R}^m):\quad\mathrm{L.\ C.}\,(\mathscr{V}_m)\to\Pi(m)-\mathrm{Ens},
\]
\textit{est une équivalence de catégories} [\textit{i. e.} $\Pi_1(\tilde{\mathscr{V}}_m)\overset{\sim}{\to}\Pi(m)$].

La théorie du degré de Brouwer fourni un homomorphisme de groupes :
\[
\deg_m:\quad\Pi(m)\to\mu_2.
\]

\textsc{Théorème} 4 (cf. [3]). --- \textit{Les conditions ci-après sont équivalentes} :

$(a)$ \textit{l'homomorphisme} $\deg_m:\Pi(m)\overset{\sim}{\to}\mu_2$ \textit{est un isomorphisme};

$(b)$ \textit{tout homéomorphisme de} $\mathbf{R}^m$ \textit{de degré} $+1$ \textit{est « stable »} (\textit{i. e. composé fini d'homéomorphismes dont chacun est l'identité sur un ouvert non vide}).

Puisque Kirby a prouvé que pour $m\ne4$ tout homéomorphisme de $\mathbf{R}^m$ de degré $+1$ est stable (\textit{cf.} [2]) on a le résultat suivant :

\textsc{Corollaire}. --- \textit{Si} $m\ne4$ \textit{l'homomorphisme} $\deg_m::\Pi(m)\overset{\sim}{\to}\mu_2$ \textit{est un isomorphisme}.

\textsc{Définition}. --- On appelle \textit{triplet de Stokes} (en dimension $m$) tout triplet $\tau=(\Omega_m,\Omega_{m-1},\delta)$ d'un faisceau séparé $\Omega_m$ sur $\mathscr{V}_m$, d'un faisceau séparé $\Omega_{m-1}$ sur $\mathscr{V}_{m-1}$ et d'un morphisme fonctoriel :
\[
\delta:\quad\Omega_m^{\bullet}\to\Omega_{m-1}\circ\partial
\]
($\partial:\mathscr{V}_m\to\mathscr{V}_{m-1}$ est le foncteur bord) qui soit un « isomorphisme local » au sens de l'axiome $(c)$.

La notion de morphisme entre triplets de Stokes en dimension $m$ étant évidente, on note par $\operatorname{St}(m)$ la catégorie des triplets de Stokes en dimension $m$.

\textsc{Théorème} 5 (cf. [3]. --- \textit{Le foncteur} :
\[
(\Omega_m,\Omega_{m-1},\delta)\to\Omega_m(\mathbf{R}^m):\operatorname{St}(m)\overset{\approx}{\to}\Pi(m)-\mathrm{Ens},
\]
\textit{est une équivalence de catégories}.

(*) Remise le 5 janvier 1981.

[1] R. \textsc{Godement}, \textit{Théorie des faisceaux}, Herman, Paris, 1958.

[2] R. \textsc{Kirby}, \textit{Ann. of Math.}, 89, 1969.

[3] M. \textsc{Wanderley}, \textit{Théorie axiomatique de l'orientation des variétés topologiques à bord} (\textit{Thèse de} 3\textsuperscript{e} \textit{cycle}, Montpellier, 1980).

[4] SGA 4, 1\textsuperscript{er} vol., \textit{Séminaire de Géométrie algébrique de l'I.H.E.S.}, dirigé par M. \textsc{Artin}, A. \textsc{Grothendieck} et J. L. \textsc{Verdier} (\textit{Lecture Notes in Math.} 289, Springer).

\textit{U.E.R. de Mathématiques, Université des Sciences et Techniques du Languedoc\unsure{.|,|;|:}}\\
\textit{Ed. N. S. das Graças, App. 202, Rua das Graças, 277, 50.000 Recife, P. E., Bresil.}

C. R., 1981, 1\textsuperscript{er} \textit{Semestre}. (T. 292, N\textsuperscript{o} \textbf{3})\hfill\textbf{Série I --- 15}

% archive sheet number: 29
\scanpage{31}
$\overline{\mathcal V}={}$ catégorie des var. top. \inserted{à bord} avec immersions ouvertes

$\mathcal V^{\deleted{\unsure{*}}}={}$ catégorie des variétés topologiques sans bord

$\partial:\overline{\mathcal V}\longrightarrow\mathcal V$ foncteur bord topologique, $\deleted{\unsure{i|X}}X\longrightarrow\deleted{\unsure{i}}\operatorname{int}(X)\inserted{=X-\partial X}\unsure{=|:}\overline{\mathcal V}\longrightarrow\mathcal V$ \margin{foncteur ``pts intérieurs}

$\overline{\mathcal V}_{\leq n},\mathcal V_{\leq n},\overline{\mathcal V}_n,\mathcal V_n$ ss-catégories des variétés de dim $\leq n$ (resp. purement de dim $n$)

Une \uline{théorie de l'orientation} des variétés topologiques consiste en

(i) Un foncteur
\[
\Omega:\mathcal V^{\unsure{o|0|\circ}}\longrightarrow(\mathrm{Ens})
\]
\margin{On note $\Omega_n=\Omega|\mathcal V_n$\\ $\Omega(X)={}$ ensemble des orientations de $X$}

(ii) Une \deleted{\unsure{application|morphisme}} élément \deleted{de}
\[
\omega^+\in\Omega(e)\qquad(e\ \text{la variété ponctuelle})
\]

(iii) \deleted{\unsure{Pour tout|Donnée}} Un hom de foncteurs sur $\overline{\mathcal V}^{o}$
\[
\Omega(\operatorname{int}X)\longrightarrow\Omega(\partial X),
\]

Ces données satisfaisant aux conditions suivantes

\uline{1}$^\circ$) $\Omega$ est un \uline{faisceau} [i.e. $\forall X\in\operatorname{Ob}\mathcal V$, la restriction \inserted{$\Omega_X$} de $\Omega$ à la catégorie des ouverts de $X$ est un faisceau \unsure{$\}$} (faisceau d'orientation de $X$)]
% slanted
\margin{\unsure{nature locale|notion locale}}
\deleted{\unsure{de plus les} $\Omega_X\to X$ sont \unsure{couvrants|ouverts|surjectifs|revêtements} (\unsure{tte variété|les variétés} \unsure{a} \unsure{local$^{\rm mt}$ une orient.|localement une orientation})}

2$^\circ$) $\operatorname{card}(\Omega(e))=2$. [L'élément de $\Omega(e)$ distinct de $\omega^+$ est noté $\omega^-$. On définit une bijection $\Omega(e)\simeq\{+1,-1\}\subset\mathbb Z$ par $\omega^+\mapsto+1$, $\omega^-\mapsto-1$ : multiplicité (égale à $+1$ ou $-1$) d'un point orienté Catégorie $\vec{\mathcal V}^{\unsure{(\Omega)}}$ des variétés \unsure{top.|top} ``orientées'' (i.e. $\Omega$-orientées)]
% slanted
\margin{hyp. de \unsure{démarrage} des \unsure{récurrences|revêtements} (\unsure{normalis$^{\rm tion}$|normalisation})}

\deleted{\unsure{$3^\circ$})} Pour tte $\vec X\in\operatorname{Ob}\vec{\mathcal V}$ et $Y$ sous-variété de codim $1$, le faisceau $\unsure{\rho|p}_{Y/X}^{\unsure{(X'',\omega)|(X,\omega)}}$ des ``rives'' de $Y$ dans $X$ (qui est un rev. étale d'ordre $2$ de \unsure{$x$|$X$}) est envoyé, grâce à la donnée (iii), dans le faisceau d'orientation $\Omega_Y$. On trouve donc, \deleted{pour tout $X\in$} $\Omega(X)\to\operatorname{Hom}(\rho_{Y/X},\Omega_Y)$, et par faisceautisation
\[
\Omega_X|Y\longrightarrow\underline{\operatorname{Hom}}(\rho_{Y/X},\Omega_Y)
\]
ou encore
\[
\Omega_X\times\rho_{Y/X}\longrightarrow\Omega_Y
\]
% archive sheet number: 30
